\begin{helper}
Let \(a,b,c\in\operatorname{Fin}m\) be pairwise distinct.  For the projected
pair rank
\[
  \rho(x,y)=\min(x,y)\,m+\max(x,y),
\]
where \(x\) and \(y\) are represented by their natural-number coordinates, one
of the three unordered pairs \(ab\), \(ac\), or \(bc\) is strictly latest:
either \(\rho(a,c)<\rho(a,b)\) and \(\rho(b,c)<\rho(a,b)\), or
\(\rho(a,b)<\rho(a,c)\) and \(\rho(b,c)<\rho(a,c)\), or
\(\rho(a,b)<\rho(b,c)\) and \(\rho(a,c)<\rho(b,c)\).
\end{helper}
\begin{proof}
Since \(a,b,c\) are pairwise distinct elements of \(\operatorname{Fin}m\),
their natural-number coordinates are pairwise distinct.  Hence exactly one of
the three coordinates is the smallest.

First observe the basic monotonicity of \(\rho\).  If the smaller coordinate
of an unordered pair \(xy\) is strictly smaller than the smaller coordinate of
another unordered pair \(zt\), then \(\rho(x,y)<\rho(z,t)\).  Indeed, the
larger coordinate of \(xy\) is still \(<m\).  Thus
\[
  \min(x,y)m+\max(x,y)<\min(x,y)m+m=(\min(x,y)+1)m.
\]
If \(\min(x,y)<\min(z,t)\), then
\(\min(x,y)+1\le \min(z,t)\), so
\[
  (\min(x,y)+1)m\le \min(z,t)m
      \le \min(z,t)m+\max(z,t)=\rho(z,t).
\]
Combining the inequalities gives \(\rho(x,y)<\rho(z,t)\).

If \(a\) has the smallest coordinate, then both pairs \(ab\) and \(ac\) have
smaller coordinate \(a\), while the pair \(bc\) has smaller coordinate
\(\min(b,c)>a\).  By the monotonicity just proved,
\(\rho(a,b)<\rho(b,c)\) and \(\rho(a,c)<\rho(b,c)\), giving the third
alternative.

If \(b\) has the smallest coordinate, the same argument shows that \(ab\) and
\(bc\) both have smaller coordinate \(b\), while \(ac\) has smaller coordinate
\(\min(a,c)>b\).  Hence \(\rho(a,b)<\rho(a,c)\) and
\(\rho(b,c)<\rho(a,c)\), giving the second alternative.

If \(c\) has the smallest coordinate, then \(ac\) and \(bc\) both have smaller
coordinate \(c\), while \(ab\) has smaller coordinate \(\min(a,b)>c\).  Hence
\(\rho(a,c)<\rho(a,b)\) and \(\rho(b,c)<\rho(a,b)\), giving the first
alternative.
\end{proof}
